%! Absolute Value Functions & Transformations — Unit 1, Lesson 1.3 — Slides (Beamer)
% Deck follows the lesson's gradual-release flow: targets → warm-up → hook → guided notes (I do /
% we do) → individual practice (you do alone) → debrief → close & homework. No group activity.
\documentclass[aspectratio=169, 11pt]{beamer}
\usepackage{algebra2-beamer}   % provides \forestheader{} and \sectionlabel[color]{}
% Coordinate grid: {xmin}{xmax}{ymin}{ymax}
\newcommand{\lgrid}[4]{%
  \draw[step=1, linegray!40, very thin] (#1,#3) grid (#2,#4);
  \draw[->, thick, charcoal] (#1,0)--(#2,0) node[below right, font=\tiny]{$x$};
  \draw[->, thick, charcoal] (0,#3)--(0,#4) node[left, font=\tiny]{$y$};}

\begin{document}

% TITLE SLIDE (CourseName is not defined in beamer, so the name is literal)
{
\setbeamercolor{background canvas}{bg=forest}
\begin{frame}[plain]
  \vspace{2.8cm}\hspace{0.55cm}%
  \begin{minipage}{0.9\textwidth}
    {\color{goldacc}\bfseries\small LESSON 1.3}\\[0.5cm]
    {\color{white}\bfseries\LARGE Absolute Value Functions \& Transformations}\\[1.2cm]
    {\color{forestmid}\small Algebra 2~$\cdot$~Unit 1: Linear Functions}
  \end{minipage}
\end{frame}
}

% LEARNING TARGETS
\begin{frame}[t, plain]
  \forestheader{Today's targets}
  \sectionlabel{What you will be able to do}
  \begin{itemize}\setlength{\itemsep}{6pt}
    \item Read the parent $f(x)=|x|$ as a \textbf{V}: vertex, axis of symmetry, minimum.
    \item Use \textbf{vertex form} $a|x-h|+k$ to find the vertex $(h,k)$ --- and say why the
          inside of the bars moves the graph \emph{backwards}.
    \item Decide from $a$: opens \textbf{up} (a minimum) or \textbf{down} (a maximum)? Narrower
          or wider than the parent?
    \item Read a V's features, build its rule from a graph, and connect the two solutions of
          $|ax+b|=c$ to where the V crosses $y=c$.
  \end{itemize}
  \vspace{4pt}
  \begin{block}{How today runs}
    Warm-up (5) $\to$ notes together (20) $\to$ \textbf{practice on your own} (15) $\to$ debrief
    (10) $\to$ close \& start the homework (10).
  \end{block}
\end{frame}

% WARM-UP
\begin{frame}[t, plain]
  \forestheader{Warm-Up: three things you already know}
  \sectionlabel{5 minutes --- alone, then check with a neighbour}
  \begin{enumerate}\setlength{\itemsep}{7pt}
    \item \textbf{Build the table} for $y=|x|$ at $x=-2,\dots,2$. Are the two halves the same?
          Where is the smallest $y$?
    \item \textbf{Moves from Lesson 1.1.} $y=x+3$ shifts which way? Is $y=2x$ steeper or flatter?
          What does a negative coefficient do?
    \item \textbf{Read the V.} Its lowest point, its fold line, and where it crosses the dashed
          line $y=3$.
  \end{enumerate}
  \vspace{6pt}
  \begin{block}{Hold on to this}
    The V crosses $y=3$ at $x=-3$ and $x=3$ --- the two answers to $|x|=3$ from last lesson.
    Keep that; we come back to it.
  \end{block}
\end{frame}

% HOOK
\begin{frame}[t, plain]
  \forestheader{Hook: the trail and the bridge}
  \sectionlabel[goldacc]{$H(x)=2|x-3|+1$ --- height above the creek, in hundreds of feet}
  \begin{columns}[T]
    \begin{column}{0.42\textwidth}
      \centering
      \begin{tikzpicture}[scale=0.45]
        \lgrid{-0.5}{6.5}{-0.5}{8}
        \draw[very thick, forest] (0,7)--(3,1)--(6,7);
        \fill[charcoal] (3,1) circle (3.5pt);
        \node[charcoal, font=\scriptsize, below right] at (3,1){$(3,1)$};
      \end{tikzpicture}
    \end{column}
    \begin{column}{0.54\textwidth}
      \begin{itemize}\setlength{\itemsep}{5pt}
        \item How high at the trailhead? At mile 3? At mile 6? \; \textbf{$7$, $1$, $7$.}
        \item Where is the trail lowest, and how low? \; \textbf{Mile 3, height 1 --- the bridge.}
      \end{itemize}
      \vspace{4pt}
      \begin{block}{The question that carries today}
        \textbf{Why can the height never drop below 1?} $|x-3|$ is a \emph{distance} --- never
        below $0$ --- so $2|x-3|+1$ is never below $1$. Every absolute value graph is this same
        \textbf{V}, moved, stretched, or flipped.
      \end{block}
    \end{column}
  \end{columns}
\end{frame}

% NOTES 1 — THE PARENT
\begin{frame}[t, plain]
  \forestheader{Notes 1: the parent $f(x)=|x|$ --- a V that turns once}
  \sectionlabel{One turning point, two mirror-image rays}
  \begin{columns}[T]
    \begin{column}{0.36\textwidth}
      \centering
      \begin{tikzpicture}[scale=0.5]
        \lgrid{-4}{4}{-1}{5}
        \draw[very thick, forest] (-4,4)--(0,0)--(4,4);
        \fill[charcoal] (0,0) circle (3.5pt);
      \end{tikzpicture}
    \end{column}
    \begin{column}{0.60\textwidth}
      \begin{itemize}\setlength{\itemsep}{4pt}
        \item \textbf{Vertex} $(0,0)$; \textbf{axis of symmetry} $x=0$; \textbf{minimum} $0$.
        \item Domain: all reals. Range: $y\ge0$ --- a distance never goes below zero.
        \item Decreasing on $(-\infty,0]$, increasing on $[0,\infty)$.
        \item Two rays pasted at the vertex: $y=x$ for $x\ge0$, $y=-x$ for $x<0$.
      \end{itemize}
    \end{column}
  \end{columns}
  \vspace{4pt}
  $x$ and $-x$ give the \emph{same} output, so the graph is a mirror image across the $y$-axis.
  That symmetry \emph{is} the axis of symmetry.
\end{frame}

% NOTES 2 — SHIFTS
\begin{frame}[t, plain]
  \forestheader{Notes 2: outside is vertical; inside is horizontal --- and backwards}
  \sectionlabel{Find the vertex by setting the inside to zero}
  \begin{columns}[T]
    \begin{column}{0.50\textwidth}
      \begin{itemize}\setlength{\itemsep}{5pt}
        \item $y=|x|+k$: up or down by $k$. \; $y=|x|-3$ has vertex $(0,-3)$.
        \item $y=|x-h|$: \textbf{right} by $h$. \; $y=|x-3|$ has vertex $(3,0)$; \; $y=|x+2|$ has
              vertex $(-2,0)$.
      \end{itemize}
    \end{column}
    \begin{column}{0.46\textwidth}
      \centering
      \begin{tikzpicture}[scale=0.36]
        \lgrid{-2}{6}{-1}{5}
        \draw[very thick, forest!40, dashed] (-1,1)--(0,0)--(4,4);
        \draw[very thick, forest] (-1,4)--(3,0)--(6,3);
        \fill[charcoal] (3,0) circle (3.5pt);
        \node[charcoal, font=\scriptsize] at (2,-2.2){$y=|x-3|$};
      \end{tikzpicture}
    \end{column}
  \end{columns}
  \vspace{4pt}
  \begin{block}{The rule that beats the sign flip}
    $|x-3|$ is smallest where $x-3=0$ --- at $x=3$. So $|x-3|$ moves the V \emph{right}, not left.
    For the trail: inside $=0$ at $x=3$, and the $+1$ lifts it, so the vertex is $(3,1)$ --- the bridge.
  \end{block}
\end{frame}

% NOTES 3 — THE NUMBER a  (the target misconception)
\begin{frame}[t, plain]
  \forestheader{Notes 3: the number $a$ --- up or down, narrow or wide}
  \sectionlabel[redacc]{Valley vs.\ ridge: same vertex, opposite job}
  \begin{columns}[T]
    \begin{column}{0.48\textwidth}
      \textbf{Valley: $H(x)=2|x-3|+1$} \\[2pt]
      \centering
      \begin{tikzpicture}[scale=0.26]
        \lgrid{-0.5}{6.5}{-0.5}{8}
        \draw[very thick, forest] (0,7)--(3,1)--(6,7);
        \fill[charcoal] (3,1) circle (4pt);
      \end{tikzpicture} \\[2pt]
      $a=2>0$: opens \textbf{up}. Vertex $(3,1)$; $1$ is a \textbf{minimum}; range $y\ge1$.
    \end{column}
    \begin{column}{0.48\textwidth}
      \textbf{Ridge: $R(x)=-2|x-3|+7$} \\[2pt]
      \centering
      \begin{tikzpicture}[scale=0.26]
        \lgrid{-0.5}{6.5}{-0.5}{8}
        \draw[very thick, redacc] (0,1)--(3,7)--(6,1);
        \fill[charcoal] (3,7) circle (4pt);
      \end{tikzpicture} \\[2pt]
      $a=-2<0$: opens \textbf{down} (a reflection). Vertex $(3,7)$; $7$ is a \textbf{maximum};
      range $y\le7$.
    \end{column}
  \end{columns}
  \vspace{3pt}
  \begin{block}{Test one point --- $x=5$ on the ridge}
    $R(5)=-2|2|+7=3$, \emph{below} $7$ --- so $7$ is the highest output. A negative $a$ flips the
    vertex's job. Width: $a$ is the right ray's slope --- $|a|>1$ narrower, $0<|a|<1$ wider.
  \end{block}
\end{frame}

% NOTES 4 — VERTEX FORM + BACK TO 1.2
\begin{frame}[t, plain]
  \forestheader{Notes 4: vertex form, all at once --- and back to Lesson 1.2}
  \sectionlabel{And your turn}
  \begin{columns}[T]
    \begin{column}{0.50\textwidth}
      $g(x)=-2|x-1|+3$:
      \begin{itemize}\setlength{\itemsep}{3pt}
        \item vertex $(1,3)$, axis $x=1$
        \item opens down, so $3$ is a \textbf{maximum}; range $y\le3$
        \item right $1$, up $3$, reflected, narrower
        \item $g(0)=-2|{-1}|+3=1$
      \end{itemize}
    \end{column}
    \begin{column}{0.46\textwidth}
      Where is the trail $500$ ft up? $H(x)=5$:
      \begin{align*}
        2|x-3|+1 &= 5 \\
        |x-3| &= 2 \\
        x &= 5 \;\text{ or }\; x = 1
      \end{align*}
      Last lesson's two solutions are the two places the V crosses $y=5$.
    \end{column}
  \end{columns}
  \vspace{4pt}
  \textbf{Guided Practice, together:} $g(x)=\tfrac12|x+4|-2$ --- every feature, then $g(0)$; then
  \emph{build} a rule from a V with vertex $(-1,4)$ through $(1,0)$: down $4$ over $2$, so $a=-2$
  and $y=-2|x+1|+4$.
\end{frame}

% INDIVIDUAL PRACTICE LAUNCH
\begin{frame}[t, plain]
  \forestheader{Individual Practice: on your own}
  \sectionlabel[goldacc]{15 minutes $\cdot$ silent $\cdot$ last page of your notes}
  \begin{enumerate}\setlength{\itemsep}{6pt}
    \item \textbf{From the rule.} $p(x)=-3|x-2|+5$: every feature, then $p(4)$.
    \item \textbf{From the graph.} Read the vertex, find $a$ from a second point, write the rule
          --- then check where it crosses $y=2$ by solving.
    \item \textbf{Same vertex, opposite job.} $y=|x+2|-3$ beside $y=-|x+2|-3$. A classmate says
          both have a minimum of $-3$. What did they miss?
  \end{enumerate}
  \vspace{5pt}
  \begin{block}{The rule}
    Set the inside to zero for the vertex; let the sign of $a$ say up or down. Stuck? Look \emph{up
    the page} at the notes --- not at your neighbour. Done early? Change problem~2's $a$ to
    $\tfrac12$: what moves, and what does not?
  \end{block}
\end{frame}

% DEBRIEF
\begin{frame}[t, plain]
  \forestheader{Debrief: correct your own page in a second colour}
  \sectionlabel[redacc]{10 minutes}
  \begin{enumerate}\setlength{\itemsep}{5pt}
    \item \textbf{1 --- inside to zero, then the sign.} $x-2=0 \Rightarrow$ vertex $(2,5)$;
          $a=-3<0$: opens down, $5$ a \emph{maximum}, range $y\le5$. $p(4)=-1$, below $5$. \checkmark
    \item \textbf{The must-land moment.} Problem~3: same vertex $(-2,-3)$, same axis --- but one
          V has $-3$ at the \emph{bottom} (range $y\ge-3$) and the other at the \emph{top}
          (range $y\le-3$). A negative $a$ flips the vertex's job.
    \item \textbf{2 --- build, then check.} Vertex $(1,-2)$; up $4$ over $2$ gives $a=2$:
          $y=2|x-1|-2$. Solve $2|x-1|-2=2$: $x=3$ or $-1$ --- the two crossings. \checkmark
    \item \textbf{Why $a$ does two jobs.} The V is two rays with slopes $a$ and $-a$: the sign
          says which way they leave the vertex, the size says how steeply.
  \end{enumerate}
\end{frame}

% CLOSE
\begin{frame}[t, plain]
  \forestheader{Close}
  \sectionlabel{What changed today}
  You walked in able to \emph{solve} $|ax+b|=c$. You walk out able to \emph{graph} $y=a|x-h|+k$
  --- read its vertex off the rule, say which way it opens and how wide, and see last lesson's two
  solutions as the two places a \textbf{V} crosses a horizontal line.
  \vspace{5pt}
  \begin{block}{Homework --- start it now}
    The \textbf{Homework} page in your packet. We start item~1 together and you keep going; whatever
    is left is due next class. Six problems plus an extension: features from three rules, a
    one-sign-apart pair, matching graphs and a table, the odd cases, a drone model, and one
    multiple-choice item.
  \end{block}
  \vspace{4pt}
  \textbf{Next:} Lesson 1.4 --- today's V is two rays pasted at the vertex; next, each ray becomes
  its own \emph{piece} of one rule, and you meet the staircase-shaped step function.
\end{frame}

\end{document}
